\clearpage
\labeledsubsubsection{Potential risks and threads extended scope}
The previous sub-sub-section focused on a \textit{plugin-system} /\\ a \underline{planned dependency-injection}.
However, the entire scope can be extended to also include \underline{unwanted/unplanned dependency-injections} and other similar attacks.

The main problem with most applications is that they use external dependencies.
Such libraries can be hijacked and modified, to then either include malicious code or return wrong results.
The linux operating system solves this issue by storing all binaries and shared libraries at a certain location and \underline{only} the \texttt{root} user has (/ should have) access to it.
This way, no user, nor application, can inject, hijack or modify anything.
Other operating systems, like Windows, do not have this level of security.
Only the system libraries are securely stored, most other libraries and binaries are usually accessible by the current user and any application that is run by them.
By utilizing and properly setting up tools like UAC\footnote{User-Access-Control} it is possible to archive a similar security level.

Another huge risk are \textit{mods} for games.
Many function on a similar base as the \textit{plugin-system}.
Thus, any mod could, potentially, be malicious.
This is especially the case when a user has to mod the game themselves, instead of letting a tool (such as a game-launcher) do this.
By giving the option for practically everyone to create a mod and upload it anywhere, for others to download and install it, the risk of someone publishing a malicious mod is high.
There is also no way of trusting user-statistics, comments and other indications on third-party sides.